\section{The Assaults of the Demons}\label{sec:assaults}

The demons fell by their own will and are fixed in it (\S\ref{sec:fall}).
Their natural gifts remain entire, and until the judgment the dark air is
given them as a place of punishment ordered to the exercise of men
(\ST{I q.~64 aa.~1, 4}; \S\ref{sec:q64}). What they do with that liberty
is the matter of the last question of the treatise, \latin{de impugnatione
Daemonum}, the assault of the demons (\ST{I q.~114 pr.}). The good angels
minister to man's salvation; the wicked, with the same natural powers over
imagination, sense, and bodies (qq.~110--111; \S\ref{sec:government}),
labour for his ruin; and God, who commands the one and permits the other,
orders both to the glory of the elect.

Scripture teaches the warfare plainly. ``For our wrestling is not against
flesh and blood; but against principalities and powers, against the rulers
of the world of this darkness, against the spirits of wickedness in the
high places'' (Eph 6:12). ``Your adversary the devil, as a roaring lion,
goeth about seeking whom he may devour. Whom resist ye, strong in faith''
(1 Pet 5:8--9). In Job the adversary asks leave and receives it within a
limit: ``all that he hath is in thy hand: only put not forth thy hand upon
his person'' (Job 1:12); ``he is in thy hand, but yet save his life'' (Job
2:6). He tempts the Lord in the desert and departs from him ``for a time''
(Luke 4:13). He desires to sift Simon ``as wheat,'' and the Lord answers,
``But I have prayed for thee, that thy faith fail not'' (Luke 22:31--32).
He ``transformeth himself into an angel of light'' (2 Cor 11:14), yet ``God
is faithful, who will not suffer you to be tempted above that which you are
able'' (1 Cor 10:13), and ``for this purpose the Son of God appeared, that
he might destroy the works of the devil'' (1 John 3:8), ``despoiling the
principalities and powers'' and ``triumphing over them in himself'' (Col
2:15; Douay throughout).

The Church confesses the same power and the same limit. The Council of
Trent, in the first canon on original sin, anathematizes whoever denies that
Adam by his transgression incurred death and with death \latin{captivitatem
sub eius potestate, `qui mortis' deinde `habuit imperium' (Hebr 2, 14), hoc
est diaboli}, captivity under the power of him who had the empire of death,
the devil (DS~1511); the decree on justification recalls that fallen men
were \latin{sub potestate diaboli ac mortis} (DS~1521). The Catechism
teaches that ``the power of Satan is, nonetheless, not infinite. He is only
a creature, powerful from the fact that he is pure spirit, but still a
creature,'' and that his action, though it may cause grave injuries, ``is
permitted by divine providence which with strength and gentleness guides
human and cosmic history'' (CCC~395). In the last petition of the Lord's
Prayer ``evil is not an abstraction, but refers to a person, Satan, the
Evil One, the angel who opposes God'' (CCC~2851). Paul~VI, in the audience
of 15 November 1972 (\S\ref{sec:faith}), called the devil the tempter
\latin{per eccellenza} who knows how to steal into us \latin{per via dei
sensi, della fantasia, della concupiscenza}, by way of the senses, the
imagination, and concupiscence, and he gave the Christian's defence in one
line: \latin{La grazia è la difesa decisiva}, grace is the decisive defence.
The rites and discipline by which the Church casts out demons are treated
in the study \emph{Catholic Exorcism: History, Discipline, and Pastoral
Practice} in this collection.

\sectionguard
\subsection{The Fathers on the Warfare with the Demons}

The Fathers wrote on the demons as pastors and as men of prayer who had
met them. Their teaching converges on three truths that the \emph{Summa}
will order: the demons tempt with a real and malicious power; they can
force no one; and their power is given, measured, and turned to good by
God.

\paragraph{Origen: the opposing powers and the freedom of the will.} The
third book of \emph{De principiis} gathers Scripture on ``the opposing
powers, or the devil himself,'' who contend with the human race: the
serpent in Eden, the evil spirit that troubled Saul, the lying spirit
before Micaiah, Satan at the right hand of Jesus the high priest in
Zechariah, the whole book of Job, the temptation of the Saviour, and Judas,
into whom ``Satan entered'' after the morsel. ``By all this, therefore,
holy Scripture teaches us that there are certain invisible enemies that
fight against us, and against whom it commands us to arm ourselves''
(\emph{De princ.} III.2.1). Origen then corrects ``the more simple among
the believers in the Lord Christ,'' who think ``that all the sins which men
have committed are caused by the persistent efforts of these opposing
powers,'' so that ``if there were no devil, no single human being would go
astray.'' ``Must we indeed suppose that the devil is the cause of our
feeling hunger or thirst? Nobody, I think, will venture to maintain that''
(III.2.1--2). From the natural appetites ``we receive certain initial
elements, and, as it were, seeds of sins,'' and it is when these are
indulged beyond measure and the first movements are not resisted that
``the hostile power, seizing the occasion of this first transgression,
incites and presses us hard in every way, seeking to extend our sins over a
wider field'' (III.2.2). Aquinas takes this argument into \ST{I q.~114
a.~3}.

Providence measures the contest. As the president of the games pairs boys
with boys and men with men, so ``each one is tempted in proportion to the
amount of his strength or power of resistance,'' and the Apostle promises
a way of escape ``that we should be able to bear it'' (III.2.3). Origen
adds a precise distinction: ``it is not the same thing to possess the power
of conquering and to be victorious,'' for the power given may be used
``either in a diligent manner, and then we prove victorious, or in a
slothful manner, and then we are defeated'' (ibid.). What the powers
suggest produces ``a (mental) commotion only, and an incitement instigating
us either to good or evil,'' and our freedom of will is ``preserved to us
in either case'' (III.2.4). No saint could sustain the whole assault of
these powers at once ``unless indeed the might of Him alone were to work in
him, who said, `Be of good cheer, I have overcome the world'\,'' (III.2.5).
Losses and calumnies are aimed beyond themselves, that ``we should be
driven either to excess of anger or sorrow, or to the last pitch of
despair,'' or to complain against God; and so ``every believer ought to
say, `Thou couldst have no power at all against Me, except it were given
thee from above'\,'' (III.2.6). All such events are ``brought about, not
indeed by God, and yet not without Him'' (III.2.7).

\paragraph{Athanasius: the discourse of Antony.} The \emph{Life of Antony}
sets at its centre the discourse in which the father of monks taught his
disciples about the demons (\emph{Vita Antonii} 16--43). Antony begins
with the Apostle: ``For we have terrible and crafty foes---the evil
spirits---and against them we wrestle,'' and ``Great is their number in the
air around us, and they are not far from us'' (21). Their nature is good:
``the demons have not been created like what we mean when we call them by
that name; for God made nothing evil, but even they have been made good.
Having fallen, however, from the heavenly wisdom, since then they have been
grovelling on earth,'' and ``out of envy of us Christians they move all
things in their desire to hinder us from entry into the heavens; in order
that we should not ascend up thither from whence they fell'' (22).

Their first weapon is the evil thought: ``But we need not fear their
suggestions, for by prayer, fasting, and faith in the Lord their attack
immediately fails.'' When suggestion fails they try terror, changing their
shapes: ``But not even then need ye fear their deceitful displays. For they
are nothing and quickly disappear, especially if a man fortify himself
beforehand with faith and the sign of the cross'' (23). Antony reads the
Leviathan of Job of the devil and his defeat: ``like a dragon he was drawn
with a hook by the Saviour, and as a beast of burden he received the halter
round his nostrils \dots\ And he was bound by the Lord as a sparrow, that we
should mock him'' (24). ``Since the Lord visited earth, the enemy is fallen
and his powers weakened'' (28). The very noise of the demons proves their
impotence, for one who has power acts at once: ``But the demons as they
have no power are like actors on the stage changing their shape and
frightening children with tumultuous apparition and various forms: from
which they ought rather to be despised as shewing their weakness.'' The
angel sent against the Assyrian needed no tumult: ``in quiet he used his
power'' (28).

Job is the proof: ``the devil was not the strong man, but God who
delivered Job to him to be tried. Certainly he had no power to do anything,
but he asked, and having received it, he hath wrought what he did''; and
the legion that begged to enter the swine shows the rest: ``But if they had
power not even against swine, much less have they any over men formed in
the image of God'' (29). What the demons fear is the Christian life:
``they fear the fasting, the sleeplessness, the prayers, the meekness, the
quietness, the contempt of money and vainglory, the humility, the love of
the poor, the alms, the freedom from anger of the ascetics, and, chief of
all, their piety towards Christ'' (30). Their predictions come from speed
and observation, as a horseman outruns a man on foot, ``for they know none
of those things which are not yet in existence; but God only is He who
knoweth all things before their birth'' (31); physicians, pilots, and
farmers forecast in the same way (33). ``they are cowards, and greatly fear
the sign of the Lord's Cross, since of a truth in it the Saviour stripped
them, and made an example of them'' (35). Antony teaches that the
visitation of the holy ones ``comes so quietly and gently that immediately
joy, gladness and courage arise in the soul'' (35), and the demons' inroad
brings confusion and fear (36). Signs are not the monk's boast: ``the
working of signs is not ours but the Saviour's work'' (38). Athanasius
relates how the demons vanished at the name of Christ (40), and how Satan
himself came to Antony's cell complaining that he was cursed without cause,
since ``I am become weak''; Antony answered, ``For the coming of Christ hath
made thee weak, and He hath cast thee down and stripped thee,'' and the
adversary, ``not being able to bear the burning'' of the Saviour's name,
vanished (41). The lesson is confidence: ``while the Lord is with us, our
foes can do us no hurt. For when they come they approach us in a form
corresponding to the state in which they discover us, and adapt their
delusions to the condition of mind in which they find us.'' So ``the enemy,
seeing Job fenced round with them, withdrew from him; but finding Judas
unguarded, him he took captive'' (42).

\paragraph{Chrysostom: the devil has no tyranny.} Chrysostom preached three
homilies against those who said that demons govern human
affairs, and on the power of man to resist the devil. ``Not so manifest is
the Sun, as the providence of God is clear. But nevertheless some dare to
say that Demons administer our affairs.'' God shows how they would
administer by letting the legion into the swine: ``Thus do Demons govern,''
and ``if God had entrusted the whole of this world to their authority, they
would have confused and disturbed everything'' (\emph{On the Power of Demons}
hom.~1.6). The second homily states the principle: ``he does not overcome by
force, nor by tyranny, nor through compulsion, nor through violence. Since
were this so, he would have destroyed all men'' (hom.~2.1). Why, then, is
he not taken out of the world? Because the contest crowns the diligent:
``if the antagonist were taken away he who overcomes is thereby injured''
(ibid.). His very name witnesses that his wickedness is by choice and not by
nature: he is called devil, ``the slanderer that is,'' from slandering Job
to God and God to Job (hom.~2.2). Paul himself used him ``as an
executioner,'' handing over the fornicator ``for the destruction of the
flesh,'' with the limit God set in Job (hom.~2.4). Chrysostom warns against
the excuse the devil loves: ``For he wishes extremely to attribute the cause
of our sins to himself, in order that we being nourished by these hopes,
and entering on all kinds of evil, may increase the chastening in our own
case, and may meet with no pardon from having transferred the cause to him.
Just as Eve met with none'' (hom.~2.5). ``The Devil is wicked; I grant this
indeed, but he is wicked for himself not towards us if we are wary''
(hom.~3.1). God has left him in the world ``in order that he might render
thee the stronger, in order that he may make the athlete more illustrious,''
since he ``not only does no harm to the wary and the heedful, but even
profits them, not owing to his own purpose (for that is wicked), but owing
to their courage who have used that wickedness aright'' (hom.~3.2).

\paragraph{Cassian: Abbot Serenus on the spirits of wickedness.} The
seventh and eighth \emph{Conferences} give the desert's teaching through
Abbot Serenus. The enemies ``oppose our progress in such a way that we can
think of them as only inciting to evil things and not forcing,'' and ``it
follows then that no one can be deceived by the devil but one who has chosen
to yield to him the consent of his own will'' (\emph{Conl.} VII.8). Spirit
can be joined to spirit in persuasion, but no spirit can be implanted in
another so as to hold it, ``for this is the true prerogative of Deity
alone'' (VII.10). Serenus holds that angels and souls have a fine body of
their own and that nothing is incorporeal but God (VII.13), a view the
\emph{Summa} does not share (\S\ref{sec:q51}). The demons learn our
thoughts from outward signs: ``Nobody doubts that unclean spirits can
influence the character of our thoughts, but this is by affecting them from
without by sensible influences \dots\ But they cannot possibly come near to
those which have not yet come forth from the inmost recesses of the soul''
(VII.15). Serenus likens them to thieves who scatter sand in a dark house
and listen for the ring of hidden metal: ``so these too, in order to explore
the treasures of our heart, scatter over us the sand of certain evil
suggestions'' (VII.16). ``not all devils can implant all the passions in
men, but \dots\ certain spirits brood over each sin'' (VII.17); with
beginners only the weaker spirits join battle, and none of the saints could
bear the malice of so many foes ``were it not that the merciful judge of our
contest, and president of the games, Christ Himself, equalized the
strength of the combatants, and repelled and checked their excessive
attacks'' (VII.20). The demons labour in the struggle and are confounded
when they fail (VII.21). From the swine that could not be entered without
leave, Serenus concludes: ``far more must we hold that they cannot of their
own free will enter into any one of men who are created in the image of
God, if they have not power to enter into dumb and unclean animals without
the permission of God'' (VII.22). In the eighth Conference Serenus teaches
that the air is full of spirits whom providence has hidden from human
sight (VIII.12), that the hostile princes of nations oppose one another and
Gabriel, as the prince of the Persians resisted him (Dan 10:13; VIII.13),
and that the evil powers bear the names of principalities and powers
because they rule nations and lesser spirits (VIII.14). The demons ``cannot
possibly find an entrance into the mind or body of anyone, nor have they the
power of overwhelming the soul of anyone, unless they have first deprived
it of all holy thoughts, and made it empty and free from spiritual
meditation'' (VIII.19).

\paragraph{Augustine: power given from above.} Augustine teaches the
limit of the demons' power in the \emph{City of God}: ``it is to be most
firmly believed that Almighty God can do whatever He pleases, whether in
punishing or favoring, and that the demons can accomplish nothing by their
natural power (for their created being is itself angelic, although made
malign by their own fault), except what He may permit, whose judgments are
often hidden, but never unrighteous.'' When they seem to transform things,
the demons ``do not create real substances, but only change the appearance
of things created by the true God so as to make them seem to be what they
are not,'' and ``I cannot therefore believe that even the body, much less
the mind, can really be changed into bestial forms and lineaments by any
reason, art, or power of the demons'' (\emph{De civ.\ Dei} XVIII.18). Yet
not all their works are appearances. When fire fell on Job's flocks and the
whirlwind struck his house, ``these were not deceitful appearances, and yet
they were the works of Satan to whom God had given this power'' (XX.19).

In \emph{De divinatione daemonum} Augustine explains how the demons
foretell. The subtlety of their aerial bodies gives them keener sense and
greater speed than ours, and their long life gives them experience
(\latin{rerum longe maior experientia}, \emph{De div.\ daem.} 3.7);
Augustine holds that the demons have such bodies, as the \emph{Summa} does
not (\S\ref{sec:q51}). Often they announce what they themselves are about
to do, for \latin{accipiunt enim saepe potestatem et morbos immittere, et
ipsum aerem vitiando morbidum reddere; et perversis atque amatoribus
terrenorum commodorum malefacta suadere}; at times they foresee from natural
signs, as a physician foresees an illness (5.9). Their foresight is far from
the height of the prophecy \latin{quam Deus per sanctos Angelos suos et
Prophetas operatur}. In their other predictions \latin{daemones plerumque
et falluntur et fallunt}: they are deceived when a command from above
overturns their plans, and they deceive from envy, \latin{qua hominum errore
laetantur} (6.10). Augustine had written that the demons sometimes learn
men's dispositions even when conceived in thought, from signs expressed in
the body (5.9). In the \emph{Retractations} he judged that sentence
\latin{rem \dots\ occultissimam audaciore asseveratione quam debui}, since
whether they know by bodily signs or by some other spiritual power
\latin{aut difficillime potest ab hominibus aut omnino non potest inveniri}
(\emph{Retract.} II.30). Aquinas cites both sentences in the article on
secret thoughts (\ST{I q.~57 a.~4}; \S\ref{sec:q57}). The whole doctrine
rests for Augustine on the will of the one who gives the power: \latin{Non
ergo eius voluntas cuius malitiae potestas in bonos datur, sed eius voluntas
a quo haec potestas datur debet nobis esse carissima}. Not the will of the
one whose malice receives power over the good, but the will of the one who
gives the power, must be dearest to us (\emph{De div.\ qq.\ 83} q.~79.5).

\paragraph{Gregory: the evil will and the just power.} Gregory reads the
first chapters of Job as the charter of the demons' power. When Satan
desires to tempt the holy man, he asks the Lord to put forth His hand,
``for the devil knows well that he is unable to do any thing of himself''
(\emph{Moralia} II.10.16). Then comes the sentence the tradition would
repeat:

\begin{quote}
But we must know that the will of Satan is always evil, but his power is
never unjust, for his will he derives from himself, but his power he
derives from God. For what he himself unrighteously desires to do, God does
not allow to be done except with justice. (\emph{Moralia} II.10.17)
\end{quote}

The evil spirit that troubled Saul is called both ``the Lord's spirit'' and
``an evil spirit,'' the Lord's ``by the concession of just power,'' evil by
the desire of an unjust will, ``so that he is not to be dreaded, who has no
power but by permission'' (II.10.17). In the Lord's words to Satan Gregory
marks ``the dispensations of heavenly pity, how He lets go our enemy, and
keeps him in; how He looses, and yet bridles him. He allows him some things
for temptation, but withholds him from others'' (II.11.19). Satan is said to
go forth from the presence of the Lord, who is everywhere, because ``being
freed from above from the pressure of an inward withholding, he went to the
execution of his desire'' (II.12.21), and he chose his hour with care, for
``the adversary does not only cast about what to do, but also when to do
it'' (II.13.22). In the homily on the angels Gregory assigns the restraint to
an order of the holy angels. The Powers are those \latin{ut eorum ditioni
virtutes adversae subjectae sint, quorum potestate refrenantur, ne corda
hominum tantum tentare praevaleant quantum volunt}: the adverse powers are
subject to them and are bridled by them, lest they prevail to tempt men's
hearts as much as they wish (\emph{Hom.\ in Ev.} 34.10).

\sectionguard
\subsection{The Assaults of the Demons (\ST{I q.~114})}\label{sec:q114}

Five articles: whether men are assailed by the demons; whether to tempt is
proper to the devil; whether all the sins of men are due to the assaults
or temptations of the demons; whether they can work real miracles to lead
men astray; and whether demons overcome by men are hindered from further
assaults (\ST{I q.~114 pr.}).

\paragraph{Men are assailed by the demons (a.~1).} The \emph{sed contra} is
Eph 6:12. Aquinas distinguishes ``the assault itself, and the ordering
thereof. The assault itself is due to the malice of the demons, who through
envy endeavor to hinder man's progress; and through pride usurp a semblance
of Divine power, by deputing certain ministers to assail man, as the angels
of God in their various offices minister to man's salvation. But the
ordering of the assault is from God, Who knows how to make orderly use of
evil by ordering it to good'' (\latin{ordo impugnationis ipsius est a Deo,
qui ordinate novit malis uti, ad bona ea ordinando}). Of the good angels,
``both their guardianship and the ordering thereof are to be referred to
God as their first author.'' The first objection argues that God sends
angels to guard man and would not send demons to destroy him. Aquinas
answers with two modes of assault. In instigating to sin the demons ``are
not sent by God to assail us, but are sometimes permitted to do so
according to God's just judgments.'' Sometimes their assault is a
punishment, and then they are sent, ``as the lying spirit was sent to punish
Achab'' (3 Kgs 22:21--22), yet ``they punish from hatred or envy; whereas
they are sent by God on account of His justice'' (ad~1). Gregory's
sentence on the evil will and the just power is the same doctrine
(\emph{Moralia} II.10.17). The second objection calls the fight unequal,
the weak and ignorant against the strong and astute. ``In order that the
conditions of the fight be not unequal, there is as regards man the
promised recompense, to be gained principally through the grace of God,
secondarily through the guardianship of the angels. Wherefore (2 Kings
6:16), Eliseus said to his servant: `Fear not, for there are more with us
than with them'\,'' (ad~2; 4 Kgs 6:16--17, where the servant's eyes are
opened on a mountain ``full of horses, and chariots of fire''). The flesh
and the world would suffice to exercise human weakness, ``but it does not
suffice for the demon's malice, which makes use of both,'' and ``by the
Divine ordinance this tends to the glory of the elect'' (ad~3).

\paragraph{To tempt is the devil's office (a.~2).} The \emph{sed contra} is
1 Thess 3:5, ``lest perhaps he that tempteth should have tempted you,'' with
the gloss, ``that is, the devil, whose office it is to tempt.'' ``To tempt
is, properly speaking, to make trial of something,'' and ``the immediate end
of every tempter is knowledge.'' The knowledge may serve a further end,
good, as when one tests a man to promote him, or bad, as when one tests him
``with the purpose of deceiving or ruining him.'' Man tempts God
sinfully when, as if uncertain, he presumes to make trial of God's power;
he tempts men sometimes to help and sometimes to harm. ``The devil,
however, always tempts in order to hurt by urging man into sin. In this
sense it is said to be his proper office to tempt,'' and a man who tempts
thus ``does this as minister of the devil.'' God tempts ``that He may
know, in the same sense as that is said to know which makes others to
know'': ``the Lord your God trieth you, that it may appear whether you love
him'' (Deut 13:3). The flesh and the world tempt ``as the instruments or
matter of temptations,'' which the devil also uses. The demons know what
happens outwardly among men, ``but the inward disposition of man God alone
knows, Who is the `weigher of spirits'\,'' (Prov 16:2); so ``the devil
tempts, in order to explore this inward disposition of man, so that he may
tempt him to that vice to which he is most prone'' (ad~2). This is
Cassian's image of the thief scattering sand to hear where the treasure
rings (\emph{Conl.} VII.16), and Antony's warning that the demons ``adapt
their delusions to the condition of mind in which they find us''
(\emph{Vita Antonii} 42). Though a demon cannot change the will, ``he can
change the inferior powers of man, in a certain degree: by which powers,
though the will cannot be forced, it can nevertheless be inclined'' (ad~3;
\ST{I q.~111 aa.~2--4}).

\paragraph{Not all sins are from the devil's temptation (a.~3).} The
objections gather strong authorities. Dionysius says ``that `the multitude
of demons is the cause of all evils, both to themselves and to
others'\,''; in Parker's version the question is how the demons,
``having fallen away from the angelic identity, as regards the desire of
the Good, have become cause of all evils both to themselves and to all the
others who are said to be corrupted'' (\emph{DN} 4.18). Damascene says
``that `all malice and all uncleanness have been devised by the
devil'\,''; in the English of his translator, ``all wickedness, then, and
all impure passions are the work of their mind'' (\emph{De fide orth.}
II.4). The \emph{sed contra} is the \emph{De ecclesiasticis dogmatibus}:
``Not all our evil thoughts are stirred up by the devil, but sometimes they
arise from the movement of our free-will.'' Aquinas distinguishes direct
and indirect cause. Indirectly, as one who dries the wood causes it to
burn, ``the devil is the cause of all our sins; because he it was who
instigated the first man to sin, from whose sin there resulted a proneness
to sin in the whole human race,'' and so the words of Damascene and
Dionysius are to be taken. Directly he is not: ``all sins are not committed
at the devil's instigation, but some are due to the free-will and the
corruption of the flesh. For, as Origen says (Peri Archon iii), even if
there were no devil, men would have the desire for food and love and such
like pleasures; with regard to which many disorders may arise unless those
desires are curbed by reason, especially if we presuppose the corruption of
our natures. Now it is in the power of the free-will to curb this appetite
and keep it in order.'' The sins that are due to the devil men commit
``through being deceived by the same blandishments as were our first
parents,'' as Isidore says. Whoever sins without the devil's instigation
still becomes his child ``in so far as he imitates him who was the first to
sin'' (ad~2; John 8:44). The comparison with the good angels fails: ``Man
can of his own accord fall into sin: but he cannot advance in merit without
the Divine assistance, which is borne to man by the ministry of the angels.
For this reason the angels take part in all our good works: whereas all our
sins are not due to the demons' instigation. Nevertheless there is no kind
of sin which is not sometimes due to the demons' suggestion'' (ad~3).

The article gathers the Fathers' teaching: Origen against those who think
no one would sin without the devil (\emph{De princ.} III.2.1--2),
Chrysostom against those who would shift the cause of their sins to him
(\emph{On the Power of Demons} hom.~2.5), and Chrysostom's reading of Eve,
who yielded ``to her own desire'' more than to the serpent's craft
(hom.~3.4). Paul~VI taught the same in 1972 with a reference to the
\emph{Summa}: \latin{Non è detto che ogni peccato sia direttamente dovuto ad
azione diabolica}; but whoever does not keep watch over himself with a
certain moral rigour exposes himself to the \latin{mysterium iniquitatis}
(audience of 15 November 1972).

\paragraph{The demons' wonders are not true miracles (a.~4).} The
\emph{sed contra} is Augustine: ``Often by means of the magic art miracles
are wrought like those which are wrought by the servants of God.'' In the
strict sense ``the demons cannot work miracles, nor can any creature, but
God alone: since in the strict sense a miracle is something done outside
the order of the entire created nature, under which order every power of a
creature is contained'' (\ST{I q.~110 a.~4}). In a wide sense a miracle is
``whatever exceeds the human power and experience,'' and so ``demons can
work miracles, that is, things which rouse man's astonishment.'' ``Although
these works of demons which appear marvelous to us are not real miracles,
they are sometimes nevertheless something real. Thus the magicians of
Pharaoh by the demons' power produced real serpents and frogs'' (Exod
7:11--12; 8:7), and the fire and wind that destroyed Job's household ``were
the work of Satan, not phantoms,'' as Augustine says (\emph{De civ.\ Dei}
XX.19). The works of Antichrist, whose coming is ``according to the working
of Satan, in all power and signs and lying wonders'' (2 Thess 2:9), are
called lying, as Augustine explains, either because he will deceive men's
senses by phantoms, ``or because, if he work real prodigies, they will lead
those into falsehood who believe in him'' (ad~1). At its locus Augustine
leaves the reason of the name to the time itself: ``Why they are called
signs and lying wonders, we shall then be more likely to know when the time
itself arrives'' (XX.19).

The second reply sets the limit. Matter does not obey the demons at will,
but they ``can employ certain seeds that exist in the elements of the
world,'' as Augustine teaches (\emph{De Trin.} III.8--9), and so whatever
natural powers can produce, the demons can produce ``by the employment of
these seeds; such as the transformation of certain things into serpents or
frogs, which can be produced by putrefaction. On the contrary, those
transformations which cannot be produced by the power of nature, cannot in
reality be effected by the operation of the demons; for instance, that the
human body be changed into the body of a beast, or that the body of a dead
man return to life.'' What seems such is ``a mere semblance of reality,''
worked from within on imagination and sense (\ST{I q.~111 aa.~3--4}) or
from without, as a demon ``can from the air form a body of any form and
shape, and assume it,'' and clothe any body in another's appearance; this
is Augustine's sentence on the phantasm presented to other men's senses
(\emph{De civ.\ Dei} XVIII.18). The third objection says that if demons can
work true wonders, miracles cannot confirm the faith. Aquinas answers with
Augustine: ``When magicians do what holy men do, they do it for a different
end and by a different right. The former do it for their own glory; the
latter, for the glory of God: the former, by certain private compacts; the
latter by the evident assistance and command of God, to Whom every creature
is subject'' (ad~3; \emph{De div.\ qq.\ 83} q.~79.4, where the Latin runs
\latin{et diverso fine et diverso iure fiunt}). \emph{De potentia} gives
the deeper reason why God grants the demons no power beyond their nature:
a miracle is a divine testimony to God's power and truth, and \latin{si
Daemonibus \dots\ aliqua potestas daretur faciendi miracula, Deus falsitatis eorum testis existeret; quod divinam
bonitatem non decet} (\emph{De pot.} q.~6 a.~5). What the demons do is
done \latin{per modum artis}, by natural powers applied, never by way of
miracle (ibid.).

\paragraph{A demon overcome does not return at once (a.~5).} The \emph{sed
contra} is Matt 4:11: ``Then the devil left him,'' Christ who had overcome
him. Aquinas records two opinions: that a demon once overcome can tempt no
man again, and that he can tempt others but not the same man. The second
``seems more probable as long as we understand it to be so for a certain
definite time: wherefore (Luke 4:13) it is written: `All temptation being
ended, the devil departed from Him for a time.'\,'' There are two reasons.
One is God's clemency, for, as Aquinas cites Chrysostom on Matthew, ``the devil does not tempt man for
just as long as he likes, but for as long as God allows; for although He
allows him to tempt for a short time, He orders him off on account of our
weakness.'' The other is the devil's cunning, as Ambrose says on Luke:
``The devil is afraid of persisting, because he shrinks from frequent
defeat.'' That he sometimes returns is plain from the Gospel: ``I will
return into my house from whence I came out'' (Matt 12:44). The objections
are easily solved. Christ overcame the tempter in the desert, and the
devil afterwards assailed him again through those who sought his death,
which Luke's ``for a time'' allows; and God's mercy in holding back the
conquered enemy is no incitement to sharper attack, since the enemy is
restrained, not rewarded. The Fathers give the same picture from their own
side. Antony teaches that ``if thus they are worsted they make an onslaught
in another manner'' (\emph{Vita Antonii} 23), and Serenus that ``when one
has been vanquished and retreated, he must make way for another spirit to
attack it still more vehemently'' (\emph{Conl.} VII.19). Gregory saw the
same rhythm in Job: God did not loose the enemy on everything at once, but
``that from one conflict first the enemy might return to the Lord defeated,
and that then he might grant him another encounter to be again worsted''
(\emph{Moralia} II.11.19).

\angeldossier{\ST{I q.~114 aa.~1--5}; parallels \emph{De malo} q.~16
aa.~8, 10--12; \emph{De potentia} q.~6 aa.~5, 10.}
 {Defined that by Adam's sin man incurred captivity under the power of the
 devil, who had the empire of death (Trent, Sess.~V can.~1, DS~1511; cf.\
 DS~1521). Church teaching that Satan's power is a creature's, not
 infinite, and permitted by providence (CCC~395), that the last petition
 of the Lord's Prayer asks deliverance from him (CCC~2850--2854), and that
 not every sin is directly due to diabolical action (Paul~VI, 15 November
 1972; a.~3). Common teaching that the demons assail men from envy and
 pride under God's ordering, permitted in tempting and sometimes sent to
 punish (a.~1; Eph 6:12; 3 Kgs 22:22; Gregory, \emph{Moralia} II.10.17),
 that to tempt to sin is the devil's proper office (a.~2; 1 Thess 3:5),
 that the devil is the indirect cause of all sins through the first sin
 but not the direct cause of all (a.~3; Origen, \emph{De princ.}
 III.2.1--2), and that the demons work wonders, sometimes real, but no
 miracle in the proper sense (a.~4; Augustine, \emph{De Trin.} III.8.13,
 \emph{De civ.\ Dei} XVIII.18, XX.19). Thomist position, as the more
 probable of two recorded opinions, that a demon overcome is kept from
 assailing the same man for a time (a.~5).}
 {Eph 6:12 (\emph{sed contra} of a.~1); 1 Thess 3:5 with the gloss (a.~2);
 the \emph{De ecclesiasticis dogmatibus} (a.~3), Augustine's
 \emph{Eighty-Three Questions} (a.~4), and Matt 4:11 (a.~5), the
 \emph{sed contra} authorities; 3 Kgs 22:21--22, 4 Kgs 6:16, Gen 22:1,
 Deut 13:3, Prov 16:2, John 8:44, 2 Thess 2:9, Mark 16:20, Luke 4:13, Matt
 12:44 (Douay). Dionysius, \emph{DN} 4.18, read in Parker; Damascene,
 \emph{De fide orth.} II.4, read at the locus; Origen, \emph{De princ.}
 III.2, read at the locus; Augustine, \emph{De civ.\ Dei} XVIII.18 and
 XX.19, \emph{De Trin.} III.8--9, \emph{De div.\ qq.\ 83} q.~79, read at
 the locus; Isidore, \emph{Sent.} II, Ambrose on Luke 4:13, and the
 Chrysostom on Matthew (the \emph{Opus imperfectum}, as the 1920 English
 edition notes), cited by Aquinas. Added witnesses read at the locus: Athanasius, \emph{Vita
 Antonii} 21--43; Chrysostom, \emph{On the Power of Demons} hom.~1--3;
 Cassian, \emph{Conl.} VII.8--22, VIII.12--19; Augustine, \emph{De div.\
 daem.} 3.7--6.10, \emph{Retract.} II.30; Gregory, \emph{Moralia}
 II.10--13, \emph{Hom.\ in Ev.} 34.10. Trent (DS~1511, 1521); CCC~395,
 2850--2854; Paul~VI, audience of 15 November 1972.}
 {``The more simple among the believers'' in Origen, who ascribe every sin
 to the opposing powers (\emph{De princ.} III.2.1), and those in
 Chrysostom who say that demons govern human affairs (hom.~1.6). On the
 demon overcome (a.~5): the opinion that he can tempt no man again, and
 the unqualified opinion that he can never tempt the same man again.
 Augustine's early sentence that the demons learn thoughts from bodily
 signs, which he later judged too confident (\emph{Retract.} II.30).}

\sectionguard
\subsection{``Deliver Us from Evil''}

The Lord's Prayer ends in the petition that answers the whole question:
``And lead us not into temptation. But deliver us from evil'' (Matt 6:13).
The Catechism reads the two petitions together. The first ``implores the
Spirit of discernment and strength'' in the battle ``between flesh and
spirit'' (CCC~2846); the second asks deliverance from Satan, ``the one who
`throws himself across' God's plan and his work of salvation accomplished
in Christ'' (CCC~2851). The victory is already won: ``Victory over the
`prince of this world' was won once for all at the Hour when Jesus freely
gave himself up to death to give us his life'' (CCC~2853). The Catechism
gives Ambrose's assurance to the newly baptized: ``One who entrusts himself
to God does not dread the devil. `If God is for us, who is against
us?'\,'' (CCC~2852, citing Ambrose, \emph{De sacramentis} 5.4.30). And when
the Church asks to be delivered from the Evil One, she prays ``to be freed
from all evils, present, past, and future, of which he is the author or
instigator'' (CCC~2854). The powers the demons keep, and the limits God
sets them, are gathered in Appendix~\ref{app:demons}.
